\section{Entorno y pruebas de viewrv}

El entorno de la aplicación es independiente del contenedor RTL. La versión
0.7.0 fija CPython 3.14.9 y Node.js 24.22.1, con dependencias en
\code{requirements-linux.lock} y \code{package-lock.json}. En Linux, desde
\code{tp3}, con esas versiones disponibles:

\begin{codebox}[language=bash]
python3.14 -m venv software/.venv
software/.venv/bin/python -m pip install --require-hashes \
  -r software/requirements-linux.lock
cd software/gui
npm ci --no-audit --no-fund
node node_modules/electron/install.js
npm run check
npm run test:core
npm run test:gui
\end{codebox}

La instalación descarga dependencias. Una vez preparadas, las pruebas utilizan
recursos locales. \code{check} comprueba formato, compila TypeScript y ejecuta
las pruebas unitarias de GUI; \code{test:core} ejecuta la suite Python de core,
CLI y protocolo; \code{test:gui} recorre las ventanas sobre la construcción
generada. Este último requiere una sesión gráfica X11 en Linux y conserva
resultados bajo \code{software/build/gui-check/}. \code{npm run check:all}
reúne los tres niveles.

Los adaptadores controlados y pseudo-terminales permiten comprobar el cliente
sin FPGA. Esas pruebas no acreditan por sí solas la GUI sobre Basys 3 ni la
distribución Windows. Los procedimientos de aplicación se detallan en el
\href{\projecturl/blob/\baserevision/tp3/software/README.md}{README de viewrv}.
